% Save as: module5-problem-set.tex
% ============================================================================
% Problem Set 5: Behavioral Economics — Deviations, Biases & Nudges
% Microeconomic Theory for Experimentalists & Applied Microeconomists
% Built from templates/problem_set_template.tex. Compile with pdflatex.
% ============================================================================
\documentclass[11pt]{article}
\usepackage{amsmath, amssymb, amsthm}
\usepackage[margin=1in]{geometry}
\usepackage{enumitem}

\newtheorem{question}{Question}

\title{Problem Set 5: Behavioral Economics --- Deviations, Biases \& Nudges\\
\large Microeconomic Theory for Experimentalists \& Applied Microeconomists}
\author{Due: \textbf{[date --- before Lecture 1 of Module 6]}}
\date{}

\begin{document}
\maketitle

\noindent\emph{Ground rules: collaboration on ideas is encouraged; write-ups
are individual. You may use AI tools for parts flagged with
$\langle$AI$\rangle$ and must attach the transcript. Show all calculations.}

\medskip

\noindent\textbf{Weights:} Part A --- 30 points. Part B --- 45 points.
Part C --- 25 points.

\bigskip

% ---------------------------------------------------------------------------
\section*{Part A: Theory --- Guided Calculations (30 points)}

\begin{question}[Present bias, sophistication, and the demand for
commitment --- 18 points]
An agent lives for three periods and must divide wealth $W = 12$ across
consumption $c_1, c_2, c_3$ with utility $\ln c$ each period. She has
$\beta$--$\delta$ preferences with $\beta = 0.7$ and $\delta = 1$: standing
in any period, utility is
\[
  \ln c_{\text{now}} \;+\; 0.7 \times \big(\text{sum of } \ln c
  \text{ in all remaining periods}\big).
\]

\begin{enumerate}[label=(\alph*)]
  \item \textbf{The exponential benchmark (3 pts).} If $\beta = 1$, argue in
        two sentences (no calculus needed --- think symmetry) that the agent
        consumes $(4, 4, 4)$.
  \item \textbf{The period-1 plan (4 pts).} With $\beta = 0.7$, the period-1
        self solves: maximize $\ln c_1 + 0.7\ln c_2 + 0.7\ln c_3$ subject to
        $c_1 + c_2 + c_3 = 12$. Show that the plan is
        $c_1 = 5,\; c_2 = c_3 = 3.5$.
        \emph{Hint: at the optimum, a dollar moved between any two periods
        must not raise utility, so $1/c_1 = 0.7/c_2 = 0.7/c_3$. Substitute
        into the budget.}
  \item \textbf{The naif's period-2 betrayal (4 pts).} Period 2 arrives with
        $12 - 5 = 7$ remaining. The period-2 self maximizes
        $\ln c_2 + 0.7 \ln c_3$ subject to $c_2 + c_3 = 7$. Show she chooses
        $c_2 \approx 4.12,\; c_3 \approx 2.88$ --- not the $(3.5, 3.5)$ the
        period-1 self planned. State in one sentence what the naive agent got
        wrong.
  \item \textbf{Commitment demand as a diagnostic (5 pts).} A bank offers a
        free account that lets the agent choose $c_1$ freely but locks the
        remainder into equal halves: $c_2 = c_3 = (12 - c_1)/2$.
        \begin{enumerate}[label=(\roman*)]
          \item Compute the period-1 utility of the \emph{sophisticated}
                agent with and without the account (with the account she
                still picks $c_1 = 5$; without it she correctly foresees
                part (c)). Show the account raises her utility.
          \item Explain in two sentences why the \emph{naive} agent's
                willingness to pay for this account is exactly zero.
        \end{enumerate}
  \item \textbf{Interpretation (2 pts).} In one paragraph: what does this
        result let a field researcher \emph{measure} by randomizing the
        \emph{price} of a commitment account? Connect to the 28\% take-up
        rate in Ashraf, Karlan \& Yin (2006).
\end{enumerate}
\end{question}

\begin{question}[Inequity aversion with numbers --- 12 points]
A responder in a \$10 ultimatum game has Fehr--Schmidt utility
\[
  U = x_{\text{self}} - \alpha \max(x_{\text{other}} - x_{\text{self}}, 0)
      - \gamma \max(x_{\text{self}} - x_{\text{other}}, 0).
\]

\begin{enumerate}[label=(\alph*)]
  \item \textbf{(4 pts)} Show that the responder accepts an offer of $s$
        dollars if and only if $s \ge 10\alpha / (1 + 2\alpha)$. Compute the
        threshold for $\alpha = 0$, $\alpha = 0.5$, and $\alpha = 1.5$.
  \item \textbf{(4 pts)} A population of responders has $\alpha$ uniformly
        distributed on $[0, 1.5]$. What fraction rejects an offer of \$2?
        Of \$3? \emph{Hint: invert the threshold formula to find the
        critical $\alpha$ for each offer.}
  \item \textbf{(4 pts)} In one paragraph: which experimental moments
        identify $\alpha$ and which identify $\gamma$? (Think about which
        game puts the decision-maker \emph{behind}, and which puts her
        \emph{ahead}.) Why can't the ultimatum game alone pin down both?
\end{enumerate}
\end{question}

% ---------------------------------------------------------------------------
\section*{Part B: Experimental Design (45 points)}

\begin{question}[Measuring the sophistication distribution, not just mean
present bias]
Part A established that commitment demand is a \emph{diagnostic for
sophistication}. Design a field experiment in a savings context whose goal
is to measure the \emph{distribution of sophistication} in the study
population --- not merely whether a commitment product raises savings.
Specify:
\begin{enumerate}[label=(\alph*)]
  \item \textbf{(6 pts)} Subject pool and recruitment, and why this
        population makes the question worth answering.
  \item \textbf{(12 pts)} Treatments. Your design must include a
        commitment-device arm with a \emph{randomized price} (possibly
        including negative prices, i.e.\ small bonuses for take-up), and you
        must explain what the take-up curve as a function of price
        identifies that a single free-offer arm cannot.
  \item \textbf{(9 pts)} Elicitation: how you measure each subject's
        \emph{predicted} future behavior as well as their actual behavior
        (the naif is defined by the gap), and why subjects have an incentive
        to report honestly.
  \item \textbf{(9 pts)} Hypotheses, and rough sample-size reasoning: anchor
        an effect-size guess to Ashraf, Karlan \& Yin (28\% take-up; savings
        up $\sim$80\%) or to your own Lab 5 simulation, and state the
        approximate $n$ per arm it implies. Simulation-based reasoning is
        welcome; formulas are not required.
  \item \textbf{(9 pts)} The biggest confound --- e.g.\ commitment demand
        driven by protecting money from household members rather than from
        one's own future self --- and your design answer to it.
\end{enumerate}
\end{question}

% ---------------------------------------------------------------------------
\section*{Part C: Silicon Pilot Interpretation $\langle$AI$\rangle$
(25 points)}

\begin{question}[The zero-switching-cost puzzle]
You run \texttt{module5-simulation-lab.py} with one modification permitted
by the CHANGE HERE blocks: the enrollment form is described as taking
\emph{zero} time --- ``a single click, right now, on the screen in front of
you.'' Theory says this should kill the default effect for \emph{every}
persona: with literally costless switching, even a procrastinator has
nothing to postpone. Yet your \texttt{exponential} (rational, disciplined)
personas still show a 25-point default effect.

\begin{enumerate}[label=(\alph*)]
  \item \textbf{(8 pts)} Give two mechanisms that could produce this, at
        least one specific to LLMs (e.g.\ the model reproducing the
        \emph{famous result} rather than responding to the decision problem;
        defaults read as advice; persona text too weak to override priors).
  \item \textbf{(9 pts)} Design the diagnostic simulation that separates
        your two mechanisms --- what do you change in the prompt or grid,
        and what result pattern favors which mechanism? Run it if API budget
        allows and report; otherwise state the predicted patterns precisely.
  \item \textbf{(8 pts)} State the design change a \emph{human-subjects}
        default study should adopt as a result of this pilot, and one
        sentence on what the pilot could \emph{not} tell you (limits of
        silicon subjects).
\end{enumerate}
Attach your AI transcripts and, if you ran code, the results CSV.
\end{question}

\end{document}
